\subsection*{1.6}\phantomsection \label{sga3:VIIB:1.6}\addcontentsline{toc}{subsection}{\SGARefLink{sga3:VIIB:1.6}{1.6}. 平展 形式簇}

为结束形式簇的一般讨论，还需简要定义在 \(k\) 上 平展 的形式簇。\footnote{\label{sga3:VIIB:nde:70}原文接着写道：“\(k\) 上的形式簇 \(X\) 称为 平展，若对角态射 \(\Delta_X:X\to X\times X\) 为局部同构；也就是说，对 \(X\) 的每个点 \(x\)，\(\Delta_X\) 诱导从 \(\mathcal O_{X\times X,\Delta_X(x)}\) 到 \(\mathcal O_{X,x}\) 的同构。由 SGA~I 容易推出，这一表述分别等价于下面两条：形式簇 \(X\) 拓扑平坦，且对每个投影到 \(s\in\operatorname{Spf}(k)\) 的点 \(x\in X\)，\(\mathcal O_{X,x}\widehat\otimes_k\kappa(s)\) 为剩余域 \(\kappa(s)\) 的有限可分扩张；或者，若 \(A\) 表示 \(X\) 的仿射代数，则对 \(k\) 的每个开理想 \(\mathfrak l\)，\(A\widehat\otimes_k(k/\mathfrak l)\) 的局部分量 \textup{（\SGARefLink{sga3:VIIB:0.1}{0.1}）} 为 \(k/\mathfrak l\) 上的有限 平展 代数。”在下文中，编者补正了上述第一条条件中遗漏的平坦性假设，并详细证明了这些条件的等价性。}

\medskip\noindent\textbf{回顾 \SGARefLink{sga3:VIIB:1.6.A}{1.6.A}。}\phantomsection \label{sga3:VIIB:1.6.A}\textit{(i)} 首先回顾 \textup{（参见 EGA \(0_{\mathrm{IV}}\)，19.10.2）}，拓扑 \(k\)-代数 \(A\) 称为\emph{在 \(k\) 上形式 平展}——相对于 \(k\)、\(A\) 各自给定的拓扑，或分别相对于离散拓扑——若对每个离散拓扑 \(k\)-代数 \(C\) \textup{（不必是 Artin 环）} 及 \(C\) 的每个幂零理想 \(J\)，每个连续 \(k\)-代数态射 \[
  A\longrightarrow C/J
\] 都唯一地提升为连续态射 \(A\to C\)，其中 \(A\) 分别赋予给定拓扑或离散拓扑。

此性质显然在基变换下保持。因此，对伪紧环的每个态射 \(k\to k'\)，\[
  A\widehat\otimes_k k'
\] 在 \(k'\) 上形式 平展。验证提升条件时也只需考虑平方为零的理想 \(J\)；参见 EGA \(\mathrm{IV}_4\)，17.1.2~(ii)。

代数 \(A\) 称为\emph{在 \(k\) 上 平展}，若它相对于离散拓扑在 \(k\) 上形式 平展，并且是有限呈示 \(k\)-代数；参见 EGA \(\mathrm{IV}_4\)，17.3.2~(ii)。下文当 \(k\) 为伪紧环、\(A\) 为预有限 \(k\)-代数时，除非另有说明，“形式 平展”均指相对于给定拓扑而言。

\smallskip\noindent\textit{(ii)} 还须回顾，若 \(F\in k[T]\) 是次数为 \(d\geq1\) 的可分首一多项式，即 \(F\) 及其导数 \(F'\) 生成的理想为整个 \(k[T]\)，则 \[
  B=k[T]/(F),
\] 在 \(k\) 上形式 平展；这里它是在 \(k\) 上秩为 \(d\) 的自由模，赋予乘积拓扑。

事实上，设 \(C\) 为离散 \(k\)-代数，从而 \(k\to C\) 的核 \(\mathfrak l\) 是 \(k\) 的开理想；设 \(J\) 为 \(C\) 的平方为零的理想，\(\phi:B\to C/J\) 为连续 \(k\)-代数态射。由于 \(B\) 是有限秩自由 \(k\)-模，\(\mathfrak lB\) 是 \(B\) 的开理想，因此 \(\phi\) 的每个提升 \(\Phi:B\to C\) 都自动连续。

设 \(t\) 为 \(T\) 在 \(B\) 中的像，任取 \(\phi(t)\) 的提升 \(u_0\in C\)。则 \(F(u_0)\in J\)，因为 \(\phi(t)\) 是 \(F\) 的根。存在 \(G,H\in k[T]\) 使得 \[
  GF+HF'=1.
\]。因此 \[
  H(u_0)F'(u_0)=1-G(u_0)F(u_0).
\]。由于 \(F(u_0)\) 平方为零，右端可逆，故 \(F'(u_0)\) 可逆。我们寻求 \(h\in J\)，使 \(u=u_0+h\) 为 \(F\) 的根。这等价于 \[
  0=F(u)=F(u_0)+F'(u_0)h,
\]，其唯一解为 \[
  h=-F'(u_0)^{-1}F(u_0)\in J.
\]。同一证明若不对态射 \(k\to C\)、\(\phi\) 和 \(\Phi\) 作连续性假设，也表明 \(B\) 是 平展 \(k\)-代数。

\smallskip\noindent\textit{(iii)} 最后，若 \[
  A=A_1\times\cdots\times A_n
\] 是有限乘积，则 \(A\) 在 \(k\) 上形式 平展，当且仅当每个 \(A_i\) 均如此。\footnote{\label{sga3:VIIB:nde:71}编者指出，EGA \(0_{\mathrm{IV}}\)，19.3.5~(v) 中给出的证明有误。}只须对 \(n=2\) 证明。设 \(e=1_{A_1}\) 为满足 \(A_1=Ae\) 和 \(A_2=A(1-e)\) 的幂等元，并设给定连续态射 \[
  \phi:A\longrightarrow C/J
\]，其中 \(J\) 平方为零。多项式 \(F=X^2-X\) 可分，因为 \[
  F'=2X-1
  \qquad\text{且}\qquad
  (F')^2-4F=1.
\]。因此 \(C/J\) 的幂等元 \(\phi(e)\) 唯一地提升为 \(C\) 的幂等元 \(f\)。故 \[
  C=Cf\oplus C(1-f),
\]，且给出提升 \(\Phi\) 等价于给出分别提升 \(\phi\) 的两个限制的态射 \[
  \Phi_1:A_1\longrightarrow Cf,
  \qquad
  \Phi_2:A_2\longrightarrow C(1-f),
\]。同样的论证表明，若 \(e\) 是 \(k\) 的幂等元且 \(A=Ae\)，则 \(A\) 在 \(k\) 上形式 平展，当且仅当它在与 \(ke\) 等同的局部分量 \(k_e\) 上形式 平展。

现设 \(X\) 为仿射代数是 \(A\) 的形式 \(k\)-簇。若 \(x\) 为 \(X\) 的点，即 \(A\) 的开极大理想 \(\mathfrak m\)，且 \(s\) 为其在 \(\operatorname{Spf}(k)\) 中的像，则用 \(k_{\mathfrak m}\) 或 \(k_s\) 表示 \(k\) 对应于 \(s\) 的局部分量。

\medskip\noindent\textbf{定义 1.6.B。}\phantomsection \label{sga3:VIIB:1.6.B}\textit{下列条件等价：}\begin{enumerate}\item[\textup{(a)}]\(A\) 在 \(k\) 上形式 平展。\item[\textup{(b)}]对每个 \(\mathfrak m\in\Upsilon(A)\)，\(A_{\mathfrak m}\) 在 \(k\) 上形式 平展，或等价地，在 \(k_{\mathfrak m}\) 上形式 平展。\item[\textup{(c)}]对 \(k\) 的每个开理想 \(\mathfrak l\)，\[
    A\widehat\otimes_k(k/\mathfrak l)
    =
    A/A\mathfrak l
  \] 在 \(k/\mathfrak l\) 上形式 平展。\item[\textup{(d)}]对每个 \(\mathfrak m\in\Upsilon(A)\) 及 \(k_{\mathfrak m}\) 的每个开理想 \(\mathfrak l\)，\(A_{\mathfrak m}/A_{\mathfrak m}\mathfrak l\) 在 \(k_{\mathfrak m}/\mathfrak l\) 上形式 平展。\end{enumerate}

若这些条件成立，则称 \(X\) \emph{在 \(k\) 上 平展}，并以 \[
  \operatorname{Vaf}^{\mathrm{\acute{e}t}}/k
\] 表示 \(\operatorname{Vaf}/k\) 中由在 \(k\) 上 平展 的形式簇构成的全子范畴。\footnote{\label{sga3:VIIB:nde:72}范畴 \(\operatorname{Vaf}^{\mathrm{\acute{e}t}}/k\) 在有限投射极限下稳定。}

事实上，设 \[
  \phi:A\longrightarrow C/J
\] 为连续 \(k\)-代数态射，其中 \(C\) 为离散 \(k\)-代数，\(J\) 平方为零。则 \(I=\operatorname{Ker}(\phi)\) 为 \(A\) 的开理想，故 \(A/I\) 是 Artin 环。因此 \(I\) 仅包含于有限多个开极大理想 \[
  \mathfrak m_1,\ldots,\mathfrak m_r.
\] 中。于是它包含 \(\mathfrak m\notin\{\mathfrak m_1,\ldots,\mathfrak m_r\}\) 时各分量 \(A_{\mathfrak m}\) 的乘积。此乘积为 \(A(1-e)\)，其中 \(e\) 为 \(A\) 中满足 \[
  Ae=\prod_{i=1}^{r}A_{\mathfrak m_i}.
\] 的幂等元。因此 \(\phi(e)=1_{C/J}\)，且给出 \(\phi\) 的连续提升等价于给出诱导态射 \[
  Ae\simeq A/A(1-e)\longrightarrow C/J.
\] 的连续提升。

另一方面，由编者注~(24)，有 \[
  A\simeq\varprojlim_{\mathfrak l}A/A\mathfrak l.
\]。由这些观察及前面的回顾，容易得到所述条件的等价性。

\medskip\noindent\textbf{定义 1.6.C。}\phantomsection \label{sga3:VIIB:1.6.C} 置 \[
  \kappa(k)
  =
  \prod_{s\in\operatorname{Spf}(k)}\kappa(s),
\]，赋予乘积拓扑。于是形式簇 \(\operatorname{Spf}(\kappa(k))\) 是各概形 \(\operatorname{Spec}\kappa(s)\)（\(s\in\operatorname{Spf}(k)\)）的直和。以 \(S_{\kappa(k)}\) 表示同样这些 \(\operatorname{Spec}\kappa(s)\) 的直和所成的概形。

对 \(k\) 上的每个形式簇 \(X\)，置 \[
  X_\kappa
  =
  X\widehat\otimes_k\kappa(k).
\]。这是经基变换得到的 \(\kappa(k)\) 上的形式簇。它与 \(X\) 有相同的点，且若 \(x\in X\) 投影到 \(s\in\operatorname{Spf}(k)\)，则 \[
  \mathcal O_{X_\kappa,x}
  =
  \mathcal O_{X,x}\widehat\otimes_k\kappa(s).
\]。由 \SGARefLink{sga3:VIIB:1.6.A}{1.6.A}~\SGARefLink{sga3:vii:paragraph:5-1:equation:i}{(i)}，基变换函子 \[
  \operatorname{Vaf}/k
  \longrightarrow
  \operatorname{Vaf}/\kappa(k)
\] 将 \(\operatorname{Vaf}^{\mathrm{\acute{e}t}}/k\) 映入 \(\operatorname{Vaf}^{\mathrm{\acute{e}t}}/\kappa(k)\)。
